\section{Erd\H{o}s Problem 43: both proposed bounds are false}

\subsection{1) Problem and source status}
Let $f(N)$ be the largest size of a Sidon subset of $[N]=\{1,\ldots,N\}$.
For Sidon sets $A,B\subset[N]$ satisfying
$(A-A)\cap(B-B)=\{0\}$, put
\[
 Q(A,B)=\binom{|A|}{2}+\binom{|B|}{2}.
\]
The questions ask whether $Q(A,B)\leq\binom{f(N)}2+O(1)$, and whether,
when $|A|=|B|$, there is a fixed $c>0$ giving the stronger asymptotic bound
\[
 Q(A,B)\leq(1-c+o(1))\binom{f(N)}2.
\]
The current page, checked on 23 September 2026, reports both answers
negative. The older local files \texttt{gpt\_pro\_5.4/43.tex} through
\texttt{43\_v9.tex} already discuss the second disproof but do not
incorporate the general solution of Problem 42 into the first question.
This note reconstructs the two known disproofs and states their
dependencies explicitly. No new discovery is claimed.

\subsection{2) The additive constant cannot exist}
Use the now-known theorem for Problem 42: for each $m\geq1$, there is
$N_0(m)$ such that every nonempty Sidon set $A\subset[N]$, $N\geq N_0(m)$,
has a Sidon companion $B\subset[N]$ of size $m$ with disjoint nonzero
differences. The adjacent \texttt{42.tex} supplies a complete deduction
from the explicitly cited compact Cayley theorem.

Define
\[
 g(N)=\max Q(A,B),
\]
where the maximum ranges over the feasible pairs. Choose $A$ with
$|A|=f(N)$; the maximum exists because $[N]$ is finite. For every
$N\geq N_0(m)$, Problem 42 gives a suitable companion of size $m$, and
therefore
\[
 g(N)-\binom{f(N)}2\geq\binom m2.
\]
Since $m$ is arbitrary, this proves the stronger conclusion
\[
 g(N)-\binom{f(N)}2\longrightarrow+\infty\qquad(N\to\infty).
\]
To check the quantifiers, given any constant $K$, choose $m$ with
$\binom m2>K$. The displayed inequality then exceeds $K$ for every
$N\geq N_0(m)$, rather than only for one finite counterexample.
This disproves the first question. It requires neither a quantitative
threshold in Problem 42 nor an asymptotic formula for $f(N)$.

\subsection{3) A self-contained form of the Bose--Chowla construction}
We next reconstruct Kevin Barreto's parity argument for the second
question, including the finite-field construction it uses.
Let $q$ be an odd prime power. In the quadratic extension
$\mathbb F_{q^2}$ choose a generator $\theta$ of its multiplicative
group, of order $v=q^2-1$. Put
\[
 S=\{s\in\{0,\ldots,v-1\}:\theta^s=\theta+t
       \hbox{ for some }t\in\mathbb F_q\}.
\]
Because $\theta\notin\mathbb F_q$, every $\theta+t$ is nonzero,
and distinct $t$ give distinct exponents. Thus $|S|=q$ and $0\notin S$.

Suppose $s_1+s_2\equiv s_3+s_4\pmod v$, and write
$\theta^{s_i}=\theta+t_i$. Multiplication and cancellation of $\theta^2$
give
\[
 (t_1+t_2-t_3-t_4)\theta+(t_1t_2-t_3t_4)=0.
\]
The elements $1,\theta$ are linearly independent over $\mathbb F_q$.
Consequently $t_1+t_2=t_3+t_4$ and $t_1t_2=t_3t_4$.
The two pairs are roots, with multiplicities, of the same monic quadratic,
so $\{t_1,t_2\}=\{t_3,t_4\}$ and
$\{s_1,s_2\}=\{s_3,s_4\}$. This proves that $S$ is Sidon modulo $v$.
Only the standard existence of finite fields and a generator of their
multiplicative groups is used in this construction.

In any abelian group a Sidon set has unique representations of nonzero
directed differences: if $s_1-s_2=s_3-s_4\neq0$, then
$s_1+s_4=s_3+s_2$, and the Sidon condition forces
$(s_1,s_2)=(s_3,s_4)$.

\subsection{4) Splitting by parity and obtaining equal sizes}
Set $N=v/2=(q^2-1)/2$. Let $E$ and $O$ be the even and odd exponent
classes of $S$, with sizes $a$ and $b$. Because $v$ is even, parity is
well-defined modulo $v$. The nonzero directed differences within $E$
and within $O$ occupy mutually distinct nonzero even residues, of which
there are exactly $N-1$. Thus
\[
 a(a-1)+b(b-1)\leq N-1=\frac{q^2-3}{2}.
\]
As $a+b=q$, this is equivalent to
\[
 (a-b)^2\leq2q-3.
\]
Put $m=\min(a,b)$, and trim the larger class so that
$E'\subset E$ and $O'\subset O$ both have size $m$. Then
\[
 m\geq\frac{q-\sqrt{2q-3}}2=\frac q2-O(\sqrt q).
\]
Define
\[
 A=\{e/2+1:e\in E'\},\qquad
 B=\{(o+1)/2:o\in O'\}.
\]
Using representatives between $1$ and $v-1$ shows that both sets lie in
$[N]$. They have equal size $m$ and are Sidon: an equality of sums lifts,
after multiplication by two and cancellation of the translation, to an
equality of sums in the modular Sidon set $S$.

If a nonzero integer $d$ were a common difference, then
\[
 e_1-e_2=2d=o_1-o_2
\]
for members of $E'$ and $O'$. This common difference is nonzero modulo
$v$, since distinct representatives lie between $1$ and $v-1$.
Uniqueness of nonzero modular directed differences would force
$(e_1,e_2)=(o_1,o_2)$, contrary to parity. Hence
$(A-A)\cap(B-B)=\{0\}$.

Since $m\leq q/2$, the construction gives
\[
 Q(A,B)=m(m-1)=\frac N2+O(N^{3/4}).
\]
In particular, $Q(A,B)\geq N/2-O(N^{3/4})$ for infinitely many
$N$, as odd primes $q$ tend to infinity.

\subsection{5) The necessary upper bound for $f(N)$}
For completeness, we prove the classical estimate sufficient to compare
the construction with $\binom{f(N)}2$. Let $C\subset[N]$ be Sidon,
$s=|C|$, and let $H\geq1$ be an integer. Define
\[
 r(x)=\#\{(c,h)\in C\times[H]:c+h=x\}.
\]
Its support has at most $N+H-1$ elements and $\sum_xr(x)=sH$, so
Cauchy--Schwarz gives
\[
 \frac{s^2H^2}{N+H-1}\leq\sum_xr(x)^2.
\]
Expanding the right side by differences, the zero difference contributes
$sH$. Each nonzero difference of $C$ has at most one representation, so
\[
 \sum_xr(x)^2
 \leq sH+2\sum_{d=1}^{H-1}(H-d)
 =sH+H(H-1).
\]
Consequently
\[
 s^2\leq(N+H-1)\left(1+\frac{s-1}{H}\right).
\]
Distinct positive differences also give $s(s-1)/2\leq N-1$, hence
$s=O(\sqrt N)$. Choosing $H=\lceil N^{3/4}\rceil$ yields
\[
 f(N)^2\leq N+O(N^{3/4}),\qquad
 \binom{f(N)}2\leq\frac N2+O(N^{3/4}).
\]
The diagonal contribution was $sH$, not $s^2H$; this is the distinction
needed for a useful convolution estimate.

For any proposed $0<c<1$, the upper bound on $f(N)$ makes
\[
 (1-c+o(1))\binom{f(N)}2
 \leq (1-c+o(1))\frac N2,
\]
whereas the constructed equal-size pairs have score
$(1+o(1))N/2$. These inequalities conflict along the infinite sequence
$N=(q^2-1)/2$. Values $c\geq1$ fail immediately as well.
This proves the second disproof without importing the full asymptotic
formula $f(N)\sim\sqrt N$.

\subsection{6) Verification and outcome}
The first argument uses the known general theorem for Problem 42 as its
external mathematical input. The second argument is written out here
from finite-field algebra, unique directed differences, parity counting,
and Cauchy--Schwarz. The two constructions serve different purposes:
the first makes the additive excess diverge, while the second rules out
a fixed proportional saving for equal-sized sets.

The accompanying \texttt{verification/solved\_42\_43\_checks.py}
constructs the finite fields explicitly for odd primes through $101$,
checks every directed modular difference, verifies the parity bound,
and checks the two final integer sets. These finite checks illustrate
the infinite-family argument; the general proof is the algebra above.

\textbf{Outcome: SOLVED, reconstruction of both known disproofs.}
The attribution and the dependency on Problem 42 are part of this
conclusion. No fresh Lean verification is asserted.

\subsection{7) Sources and attribution}
\begin{itemize}
\item T. F. Bloom, Problem 43 and its discussion,
\texttt{https://www.erdosproblems.com/43} and
\texttt{https://www.erdosproblems.com/forum/discuss/43}, checked
23 September 2026. Kevin Barreto's comment of 19 December 2025 gives
the parity construction. The implication from Problem 42 is also
recorded on the current problem page.
\item Problem 42 and the compact Cayley theorem, cited and precisely
stated in the adjacent \texttt{42.tex}; the formal source is
\texttt{https://github.com/Shashi456/erdos-formalizations/blob/main/Erdos/P42/CompactCayley/Proof.lean}.
\item The Bose--Chowla construction is proved in Section 3 rather than
used as an unexplained assumption. Its use for this problem follows
Barreto's cited argument.
\end{itemize}

This measures coverage of both original questions by the known proofs,
with the stated theorem for Problem 42 accepted. It is not a novelty
claim or a measure of independent formal verification.
\subsection{8) COMPLETION ESTIMATE}
\textbf{COMPLETION: 100\%}
